\clearpage
\VerdierPhysicalPageAnchor{0234}%
\begingroup\small\setlength{\parskip}{0pt}\setlength{\emergencystretch}{1.5em}
\setlength{\abovedisplayskip}{4pt plus 2pt minus 2pt}%
\setlength{\belowdisplayskip}{4pt plus 2pt minus 2pt}%

\VerdierSemanticAnchor{ch4d-4-1-9-cont}
\noindent Désignons par
$\bigl({}^{1}\!Y(p,q),{}^{1}\!\delta(p,q,r)\bigr)$ (resp.\ 
$\bigl({}^{2}\!Y(p,q),{}^{2}\!\delta(p,q,r)\bigr)$) l'objet spectral, à
valeurs dans $K(\mathcal A)$, défini par le complexe $X$, filtré par son
premier (resp.\ deuxième) degré (4.1.7). Il résulte immédiatement des
définitions qu'on a les relations :

\VerdierSemanticAnchor{ch4d-4-1-9-1}
\begin{equation}
\left\{
\begin{aligned}
{}^{1}\!Y(q-1,q)&=X^{-q,\bullet}[q],\\
{}^{2}\!Y(q-1,q)&=X^{\bullet,-q}[q],\\
{}^{1}\!\delta(q-1,q,q+1)&=-\bigl(d_1^{-q-1,\bullet}\bigr)[q+1],\\
{}^{2}\!\delta(q-1,q,q+1)&=-\bigl(d_2^{\bullet,-q-1}\bigr)[q+1].
\end{aligned}
\right.
\tag{4.1.9.1}
\end{equation}
Ces formules seront utilisées au numéro (4.6).

\VerdierSemanticAnchor{ch4d-4-2}
\medskip\noindent\textbf{4.2. Objet spectral à valeurs dans $D(\mathcal A)$
associé à un objet filtré de $\operatorname{comp}(\mathcal A)$.}

\VerdierSemanticAnchor{ch4d-4-2-scope}
\medskip\noindent\emph{Dans ce numéro et les numéros suivants de ce
paragraphe, $\mathcal A$ désigne une catégorie abélienne.}

\VerdierSemanticAnchor{ch4d-4-2-1}
\medskip\noindent\textbf{4.2.1.}\quad
Soient $X$ un objet de $\operatorname{comp}(\mathcal A)$ et $X(q)$,
$q\in\mathbb Z$, une filtration croissante de type
$\widetilde{\mathbb Z}$ du complexe $X$ :

\VerdierSemanticAnchor{ch4d-4-2-1-1}
\begin{equation}
0=X(-\infty)\subset\cdots\subset X(q)\subset X(q+1)
\subset\cdots\subset X(+\infty)=X.
\tag{4.2.1.1}
\end{equation}

\VerdierSemanticAnchor{ch4d-4-2-1-associated-object}
On en déduit, par la construction rappelée en (4.1.1), un objet spectral
$\bigl(\mathcal X(p,q),\delta(p,q,r)\bigr)$ à valeurs dans
$K(\mathcal A)$, d'où, en appliquant le foncteur canonique
$K(\mathcal A)\longrightarrow D(\mathcal A)$, un objet spectral à valeurs
dans $D(\mathcal A)$ encore noté, abusivement,
$\bigl(\mathcal X(p,q),\delta(p,q,r)\bigr)$. On se propose de décrire un
objet spectral à valeurs dans $D(\mathcal A)$ isomorphe à ce dernier objet
spectral.

\VerdierSemanticAnchor{ch4d-4-2-2}
\medskip\noindent\textbf{4.2.2.}\quad
Soit donc
$0=X(-\infty)\subset\cdots\subset X(q)\subset\cdots\subset
X(+\infty)=X$, un objet filtré de $\operatorname{comp}(\mathcal A)$.
Posons :

\VerdierSemanticAnchor{ch4d-4-2-2-1}
\begin{equation}
X(p,q)=X(q)/X(p),
\tag{4.2.2.1}
\end{equation}
pour tout élément $(p,q)$ de
$\mathcal{F}\ell_1(\widetilde{\mathbb Z})$. (Le quotient est pris
évidemment ici au sens de la catégorie abélienne
$\operatorname{comp}(\mathcal A)$.) Désignons par :

\VerdierSemanticAnchor{ch4d-4-2-2-2}
\begin{equation}
\varphi\bigl((p,q),(p',q')\bigr)\colon
X(p,q)\longrightarrow X(p',q'),
\tag{4.2.2.2}
\end{equation}
\endgroup
